% !TeX root = ../document_maitre.tex
\label{chap:défis-patrimoniaux}

Aujourd'hui, les \gls{sgc} se basent sur une architecture logicielle appliquant les principes décrits précédemment. Ils sont donc assez complexes, disposant de multiples entrées de données, fonctionnalités et manières d'accomplir des tâches précises. 

Au \coeur de ces logiciels se situe la notice, c'est-à-dire l'unité contenant toutes les informations en rapport avec un objet ou une ressource disponible dans l'application. Elle se base sur les métadonnées, ces données générées par les professionnels afin de décrire et contextualiser des objets patrimoniaux. 

Le monde patrimonial a développé une manière bien particulière de créer et gérer cette métadonnée. Il s'agit de tout un univers gravitant autour de principes, standards et modèles conceptuels de gestion. Des notions que nous avons en partie définies, mais qu'il convient d'examiner plus en détail.

\subsection{La métadonnée dans le monde patrimonial}

Le patrimoine s'est toujours caractérisé par la diversité des objets conservés et des pratiques professionnelles autour de leur conservation. Au cours de leur existences, chaque institutions à développés ses propres méthodes et vocabulaires de descriptions rendant les données hététogènes et inégale d'un établissement à l'autre.

L'informatique donne les moyens de créer un réseau d'informations patrimoniales. Cependant les différences de gestion et de descriptions constituaient des freins trop importants. 

Pour cette raison, il y a eu des efforts internationaux d'homogénéisation des pratiques de descriptions et de gestion. Cela afin de mettre en place et respecter un but unique : \gls{intero}.\footcite{jonkerQuestceQueLinteroperabilite2026} Il s'agit de permettre à des systèmes informatiques de s'échanger des données, et de les interpréter de la même manière. Cela sous-entend que les données doivent pouvoir être identiques et stockées dans un schéma fixe lisible par tout systèmes.\newline

Selon ce principe, plusieurs outils ont été mis en place afin d'homogénéiser un maximum les métadonnées. D'abord, des référentiels de métadonnées afin d'utiliser les mêmes éléments pour parler de la même entité. Aussi bien à l'échelle du vocabulaire, grâce aux \gls{thesaurus}, c'est-à-dire des classements de mots en listes hiérarchique, que de référentiels plus larges comme les autorités. 

Ensuite, des normes descriptives telles que l'ISAD(G) pour les archives ou la norme MPEG-7 pour les documents multimédias, comme les \glspl{imgani}. Ces normes ont pour objectif d'établir : ce qu'il faut décrire, où trouver l'information et où la restituer dans un schéma descriptif. Si les référentiels permettent de normaliser la terminologie des métadonnées, les normes descriptives visent à normaliser leur structure. Elles ont été rapidement traduites dans des standards d'encodage imposant le respect de la hiérarchie informationnelle de la norme. C'est le cas du standard EAD pour les archives, selon la norme ISAD(G), via l'utilisation du langage XML dont on peut contraindre la structuration. 

Ces éléments sont valables pour toutes les typologies patrimoniales conservées. Ainsi, elles possèdent parfois leurs propres référentiels, normes et standards d'encodage. Cela a donc homogénéisé les pratiques descriptives, mais a aussi amené à plus de distinction entre les typologies patrimoniales.\newline

L'apparition de séparations plus nette entre les typologies et des standards descriptifs n'a constitué qu'une première étape. Dans le même temps, un constat de changement des attentes du public et des besoins de gestion des collections patrimoniales conduit à une nouvelle évolution. 

C'est l'apparition de modèles conceptuels de gestion des collections. Ils ont pour objectif d'apporter de nouvelles conceptions de gestion afin de mieux répondre aux enjeux de conservation et aux attentes du public. Nous avons déjà cité le modèle CIDOC-CRM du \gls{CIDOC} et le FRBR de l'\gls{IFLA}. Nous explorerons d'ailleurs ce dernier plus en profondeur par la suite\footnote{\fullref{FRBR-WEMI}}

Il est important de retenir qu'aujourd'hui ces modèles permettent de créer une stabilité de gestion des collections. Cela crée plus de liens entre les institutions ainsi qu'une stabilité des éléments nécessaires pour retrouver les bonnes informations.\newline

Les efforts de normalisation à différentes échelles sont aujourd'hui plus que nécessaires. L'explosion des moyens informatiques et des créations du patrimoine numérique ont amplifié la masse de métadonnées produites. C'est le phénomène de \gls{bgdt}.

Aujourd'hui, les institutions conservatrices gèrent des quantités impressionnantes de métadonnées descriptives. Les enjeux volumétriques rendent plus difficile l'application de ces standards et normes.

Ensuite par la diversification des formats. Cette massification des métadonnées a demandé de nouveaux moyens de gestions et de tansmissions. Les fichiers XML étant assez lourds et verbeux, d'autres formats de fichier sont apparus : c'est le cas du \J. Sans redéfinir les normes, son apparition impacte les formats d'échange de données entre institutions. Par exemple, les sites internet institutionnels se dotent de systèmes d'échange de données, souvent basés sur le format \J : ce sont les APIs. Cela implique donc de pouvoir transformer, par exemple, un encodage EAD en un fichier \J, et inversemment.

\subsection{La place des \glspl{sgc}}

L'organisme du \textit{Collection Trust} impose le respect de ces principes concernant aux \gls{sgc}. Leur mission est finalement de constituer un terrain d'entente entre les besoins institutionnels et les efforts de normalisation. À ce titre, le logiciel SkinSoft respecte chaque norme, standard et modèle conceptuel attribués à chaque typologie. Tout en étant directement connectés avec des référentiels externes.\newline

Cependant, ce n'est pas choses aisée. En effet, ces efforts de normalisation sont aujourd'hui en pleine application. Au moment de l'essor de l'informatisation, les institutions ont souvent constitué leurs propres bases de données selon leurs propres conceptions de gestion de la métadonnée. Ces bases dites \guillemet{maisons} sont relativement fidèles aux normes. La structuration de la donnée leurs est propre : les métadonnées sont saisies selon des standards internes et l'architecture de la base de données est souvent unique.

C'est alors un premier rôle attribué aux \gls{sgc} que de permettre cette transition d'une base de données \guillemet{maison} vers un système normalisé. Il faut pouvoir assurer la bonne transition de la donnée et des pratiques de saisie. Cela implique une notion de migration de données et d'accompagnement au changement.

De plus, les \glspl{sgc} sont en première ligne face aux défis de la \gls{bgdt}. Tout en respectant ces normes, standards et modèles conceptuels, il faut pouvoir gérer de manière homogène cette masse de données. Traiter ce volume croissant et évolutif constitue un défi d'infrastructure, mais aussi un enjeu face à cette homogénéisation. 

Cette massification des données et cette mise en danger de l'homogénéisation des métadonnées impose à ces logiciels de mettre en place des moyens suffisants pour conserver et perpétuer ce principe. Ainsi, ce logiciel ne se limite plis seulement à un outil de gestion avec une dimension stratégique mais devient un moyen d'application de l'\gls{intero} et de l'homogénéité des métadonnées.

En conclusion, les \glspl{sgc} ne plus de simples logiciels d'applications de normes. Ils sont de véritables moteurs et appuis de l'évolution du patrimoine dans sa dimension numérique. 

Aujourd'hui, ils ont pour enjeu de concilier les volontés d'homogénéisation à la réalité d'hétérogénéité massive. Cela les amène à côtoyer leurs limites et donc à se diriger vers des solutions d'\gls{IA} pour faciliter cette mission par des processus automatisés ou semi-automatisés. 